\documentclass[12pt, letterpaper]{article}

\usepackage[margin=1in]{geometry}
\usepackage{setspace}
\usepackage{parskip}
\usepackage{enumitem}
\usepackage{titlesec}
\usepackage{hyperref}
\usepackage{microtype}
\usepackage{lmodern}
\usepackage[T1]{fontenc}
\usepackage{ulem}   % for \uline
\normalem            % keep \emph as italics

\hypersetup{colorlinks=true, linkcolor=blue, urlcolor=blue}

\titleformat{\section}{\large\bfseries}{}{0em}{}
\titlespacing*{\section}{0pt}{1.2em}{0.3em}

\setstretch{1.1}
\setlength{\parskip}{0.6em}

\begin{document}

\begin{center}
    {\LARGE\bfseries Innovation Statement}\\
    \vspace{0.3em}
    {\normalsize Shrinithi Natarajan}
\end{center}

\vspace{0.5em}
\hrule
\vspace{0.8em}

\section*{B. INNOVATION}

Drug resistance accounts for the overwhelming majority of cancer deaths in patients receiving
targeted therapy \cite{liu2025multiomics, mao2025ai}. Contemporary approaches such as DRN-CDR
\cite{saranya2024drncdr} and MSDRP \cite{zhao2023msdrp} train deep neural networks on multi-omics
feature vectors, returning a probability score whose value reflects weighted correlations learned
from training data. When features conflict, the
network resolves the tension implicitly through learned weights, with no mechanism to surface,
inspect, or justify the resolution. Three recent developments make a principled alternative
viable now: (1) the Model Context Protocol (MCP), standardised in 2024, enables LLM agents to
query structured databases through typed, schema-validated tool calls rather than free text;
(2) leading LLMs now follow multi-step tool-use protocols reliably enough for structured
reasoning chains \cite{flotho2026mcpmed, widjaja2025bioinformcp}; and (3) GDSC2 and CCLE/DepMap now provide matched multi-omics profiles at
the scale and quality needed to evaluate multi-modal reasoning. Published tumour board frameworks
document a hierarchy of evidence in treatment decisions --- DNA-level findings designated as
primary determinants when they conflict with expression data \cite{evomdt2026, han2024consensus,
lorenzo2026oncotimia} --- whether encoding that hierarchy computationally improves fidelity and
accountability is the central question this proposal addresses. Our framework introduces five
innovations that have not appeared individually or in combination in any published work.

\uline{A formalised Biological Hierarchy of Truth governing inter-agent conflict resolution.}
Multi-agent debate systems in oncology \cite{evomdt2026, han2024consensus} have demonstrated that
structured disagreement between specialised agents improves clinical decision-making over
single-agent systems. However, none encodes the mechanism of resolution as explicit, testable
logic --- conflicts are settled by LLM negotiation, which is unauditable and biologically
opaque. \textbf{We encode a five-tier axiom hierarchy (T1 Structural Primacy $\to$ T2
Transcriptional Gate $\to$ T3 Pathway Bypass $\to$ T4 Pharmacological Prior $\to$ T5
Statistical Consensus) as executable rules, making the resolution process fully reproducible
and biologically grounded.} Where a probabilistic model returns \texttt{P(SENSITIVE) = 0.61} with no further account,
our framework surfaces the same conflict explicitly: Genomics Agent returns SENSITIVE (BRAF
V600E, T1), Transcriptomics Agent returns RESISTANT (BRAF silenced, T2), T1 overrides T2,
final verdict SENSITIVE, dissent logged. The resolution is inspectable, reproducible, and anchored to a biological rationale rather
than a statistical weight. A 15\% confidence decay applied per verdict flip further penalises
sycophantic convergence, ensuring agents cannot agree without a higher-tier axiom justification.
To our knowledge, this is the first formalisation of inter-modality conflict resolution in
oncology AI as executable axiom logic.

\uline{The Model Context Protocol as a structural anti-hallucination layer.}
Medical hallucination is a well-documented failure mode of LLM-based clinical tools
\cite{kim2025hallucination, nishisako2025hallucination}, and existing multi-agent oncology
systems rely on model self-restraint to prevent it. MCP --- newly standardised and now
demonstrated in bioinformatics contexts \cite{flotho2026mcpmed, widjaja2025bioinformcp} ---
provides the enabling technology: \textbf{by routing every agent data request through
MCP-validated tool calls returning typed JSON, we convert hallucination from a probabilistic
risk into a detectable protocol violation.} An agent cannot assert a mutation the genomics
server did not return; the guarantee is architectural, not probabilistic.

\uline{A debate trace and dissent log as a clinical accountability artefact.}
Every prediction is accompanied by a structured record of which agents agreed, which were
overridden, by which axiom tier, and at which round --- with overridden agents retaining their
dissent in the final output, mirroring the minority opinions preserved in real tumour board
decisions. \textbf{This debate trace has no equivalent in any published AI oncology system:
rather than a confidence score, a clinician receives a citable argument --- which data layer
drove the verdict, which agent dissented, and on what biological grounds --- enabling oversight
and post-hoc audit of every recommendation \cite{evomdt2026, lorenzo2026oncotimia}.}

\uline{A novel synthesis of multi-agent debate, structured tool access, and biological axiom logic.}
Prior work combines at most two of these elements: EvoMDT \cite{evomdt2026} uses debate
without axiom logic; BioAgents \cite{mehandru2025bioagents} uses tool-anchored reasoning
without conflict resolution; DRN-CDR \cite{saranya2024drncdr} integrates multi-omics without
transparency. \textbf{This framework occupies an unoccupied position in that design space ---
structured debate, MCP-enforced data access, and an explicit biological axiom hierarchy
simultaneously --- and in doing so creates a class of system that did not previously exist.}

\uline{Purpose-built infrastructure unavailable in any existing system.}
The framework has been developed alongside unique evaluation resources: four MCP servers
wrapping GDSC2, CCLE/DepMap, KEGG, and ChEMBL with typed, schema-validated APIs that
structurally prevent agent hallucination; 40 GDSC2 cases stratified across all five axiom tiers
with mechanistic annotations designed to stress-test each tier independently; and an
LLM-agnostic client supporting Gemini Flash, Qwen, Gemma, and Llama-3.3 under a unified
interface with SQLite-backed response caching. \textbf{This infrastructure enables the first
direct comparison of free large language models on a structured multi-omics reasoning task,
allowing the contribution of the architecture to be separated from the contribution of any
specific model.}

Taken together, these five innovations shift the evaluation target for oncology AI from
predictive accuracy alone toward \textit{reasoning quality} --- hallucination rate, biological
faithfulness of debate traces, and axiom invocation patterns as first-class scientific outcomes.
If the proposed aims are achieved, the immediate result is a validated, open-source framework
that a clinical research team could deploy to generate auditable treatment sensitivity predictions
alongside their existing genomic profiling pipeline --- replacing a probability score with a
structured biological argument. The longer-term result is a generalizable architecture for any
domain requiring principled reconciliation of conflicting multi-modal evidence.

\vspace{0.8em}
\hrule
\vspace{0.5em}

\section*{Reflection}

Reading the funded examples made clear that the five-question structure from the module guide
is a \textit{planning} tool, not a formatting template. Funded innovation sections read as a single
continuous argument: gap $\to$ novel contributions $\to$ implications. The underlined topic
phrase per paragraph (as in the TCGA immune cell example) lets a reviewer scan the novelties
quickly without breaking the prose into disconnected subsections. The most important revision was
leading with the gap rather than with enabling advances --- reviewers need to believe the problem
is unsolved before they will care about the solution. The BRAF V600E clinical vignette was
folded into the T2 axiom description as a concrete illustration rather than a standalone aside.

\section*{Feedback}

\textit{[To be completed after peer or supervisor review: who gave feedback, how it was
solicited, and what changes were made in response.]}

\vspace{0.8em}
\hrule
\vspace{0.5em}

{\small
\begin{thebibliography}{99}
\setlength{\itemsep}{0pt}

\bibitem{liu2025multiomics} Liu X et al. Multi-omics advances in understanding cancer drug resistance. \textit{Biomed Technol}. 2025. doi:10.1016/j.bmt.2025.100119

\bibitem{mao2025ai} Mao Y et al. Emerging AI-driven precision therapies in tumor drug resistance. \textit{Mol Cancer}. 2025. doi:10.1186/s12943-025-02321-x

\bibitem{evomdt2026} Liu Q et al. EvoMDT: self-evolving multi-agent system for clinical decision-making. \textit{npj Digit Med}. 2026. doi:10.1038/s41746-025-02304-8

\bibitem{han2024consensus} Han X et al. Multi-Agent Medical Decision Consensus Matrix for Oncology MDT. \textit{arXiv}:2512.14321. 2024.

\bibitem{kim2025hallucination} Kim Y et al. Medical Hallucinations in Foundation Models. \textit{arXiv}:2503.05777. 2025.

\bibitem{nishisako2025hallucination} Nishisako S et al. Reducing Hallucinations in Generative AI for Cancer Information. \textit{JMIR Cancer}. 2025. doi:10.2196/70176

\bibitem{flotho2026mcpmed} Flotho M et al. MCPmed: MCP-enabled bioinformatics for LLM-driven discovery. \textit{Brief Bioinform}. 2026. doi:10.1093/bib/bbag076

\bibitem{widjaja2025bioinformcp} Widjaja F et al. BioinfoMCP: MCP Interfaces in Agentic Bioinformatics. \textit{arXiv}:2510.02139. 2025.

\bibitem{saranya2024drncdr} Saranya KR, Vimina ER. DRN-CDR: cancer drug response prediction using multi-omics. \textit{Comput Biol Chem}. 2024. doi:10.1016/j.compbiolchem.2024.108175

\bibitem{zhao2023msdrp} Zhao H et al. MSDRP: deep learning model for drug response prediction using multi-omics. \textit{Bioinformatics}. 2023.

\bibitem{mehandru2025bioagents} Mehandru N et al. BioAgents: multi-agent systems for bioinformatics. \textit{Sci Rep}. 2025. doi:10.1038/s41598-025-25919-z

\bibitem{lorenzo2026oncotimia} Lorenzo L et al. Evaluation of Oncotimia: LLM system for tumour boards. \textit{arXiv}:2601.19899. 2026.

\end{thebibliography}
}

\end{document}
